\chapter{Dominance}

\section{Notation Importante}
Cette notation est fondamentale et sera utilisée tout au long du cours. Elle peut être déroutante au premier abord.

\begin{itemize}
    \item $X_{-i} = \prod_{j \in N, j \ne i} X_j$ désigne l'ensemble des profils d'actions de tous les joueurs \textbf{sauf} le joueur $i$.
    \item $x_{-i}$ désigne un élément de $X_{-i}$ (un vecteur d'actions des adversaires).
    \item $u_i(x_i, x_{-i})$ est une notation abusive pour $u_i(x)$. Elle met en évidence que le gain dépend du choix $x_i$ du joueur et des choix $x_{-i}$ des autres.
\end{itemize}

\begin{exemple}[Exemple à 3 joueurs]
Soit $N = \{1, 2, 3\}$. Pour le joueur $i=2$ :
\begin{itemize}
    \item Les adversaires sont $1$ et $3$.
    \item Le profil des adversaires est $x_{-2} = (x_1, x_3)$.
    \item L'ensemble de ces profils est $X_{-2} = X_1 \times X_3$.
    \item La fonction de gain s'écrit $u_2(x) = u_2(x_2, x_{-2}) = u_2(x_2, (x_1, x_3))$.
\end{itemize}
\end{exemple}

\section{Action Dominante et Équilibre}

\subsection{Définitions}
\begin{definition}
Une action $x_i \in X_i$ est \textbf{dominante} si elle rapporte toujours plus (ou autant) que toute autre action, quels que soient les choix des autres :
$$\forall d_{-i} \in X_{-i}, \forall d_i \in X_i, u_i(x_i, d_{-i}) \ge u_i(d_i, d_{-i})$$
\end{definition}

\begin{definition}
Un profil $x = (x_i)_{i \in N}$ est un \textbf{équilibre en stratégies dominantes} si pour chaque $i$, $x_i$ est dominante.
\end{definition}

\paragraph{Rationalité et Défaillance Collective :}
L'exemple ci-dessous montre qu'être rationnel individuellement (choisir sa stratégie dominante) peut mener à une catastrophe collective. C'est le paradoxe central du Dilemme du Prisonnier.

\begin{center}
\begin{tabular}{c|c|c}
P1 \textbackslash P2 & a & b \\ \hline
A & (1 000 000, 1 000 000) & (0, 1 000 001) \\ \hline
B & (1 000 001, 0) & (100, 100)
\end{tabular}
\end{center}
Ici, jouer B est dominant pour les deux, menant à (0.1, 0.1), alors que la coopération (A, A) donnait 1 Million chacun.

\begin{remarque}[Réflexion]
Avez-vous en tête des exemples dans la société où jouer une action dominante pour tout le monde mène à un mauvais résultat ? (Pollution, course aux armements, etc.)
\end{remarque}

\section{Connaissance Commune (Common Knowledge)}
\begin{definition}
Un énoncé $S$ est \textbf{connaissance commune} si :
\begin{itemize}
    \item Chacun sait que $S$ est vrai.
    \item Chacun sait que les autres savent que $S$ est vrai.
    \item Chacun sait que les autres savent que les autres savent... à l'infini.
\end{itemize}
\end{definition}
\textbf{Nuance :} Il y a une différence fondamentale entre "Je sais S" et "S est connaissance commune". Votre comportement changera radicalement selon que vous pensez que l'autre est au courant ou non (exemple de l'attaque coordonnée ou de l'e-mail).

\section{Élimination des Stratégies Strictement Dominées (IESDS)}

\subsection{Action Strictement Dominée}
\begin{definition}
$x_i$ est \textbf{strictement dominée} par $y_i$ si :
$$\forall d_{-i} \in X_{-i}, u_i(x_i, d_{-i}) < u_i(y_i, d_{-i})$$
Un agent rationnel ne jouera jamais une action strictement dominée.
\end{definition}

\subsection{Processus IESDA (Formalisation)}
\begin{itemize}
    \item On part du jeu initial $G = G(0)$ avec les ensembles d'actions $X_i(0)$.
    \item \textbf{Étape 1 :} On élimine toutes les actions strictement dominées dans $G(0)$ pour obtenir les ensembles $X_i(1)$ et le jeu réduit $G(1)$.
    \item \textbf{Récurrence :} À l'étape $k$, on élimine les actions strictement dominées dans le jeu $G(k-1)$ pour obtenir $G(k)$.
    \item \textbf{Résultat :} L'ensemble des profils survivants est l'intersection de toutes les étapes :
    $$E^{\infty} = \bigcap_{n \ge 0} X(n)$$
\end{itemize}

\begin{corollaire}[Lien Nash-IESDA]
L'ensemble des équilibres de Nash est toujours inclus dans l'ensemble des profils survivant à l'IESDA. Si l'IESDA réduit le jeu à un seul profil, c'est l'unique équilibre de Nash.
\end{corollaire}

\section{Exercices et Méthodologie}

\subsection{Méthode de résolution (Type Examen)}
Pour résoudre un jeu par IESDA, il faut écrire explicitement les étapes :
\begin{enumerate}
    \item "Dans le jeu initial, pour le Joueur $i$, l'action $A$ domine strictement $B$ car $u_i(A, \cdot) > u_i(B, \cdot)$."
    \item "On supprime l'action $B$."
    \item "Dans le jeu réduit, pour le Joueur $j$, l'action..."
\end{enumerate}

\subsection{Exercices d'Application}

\begin{exemple}[Exercice d'application (Fonctions continues)]
Soit un jeu à deux joueurs avec les gains :
\[ u_1(x, y) = (-x^2 + x + 1)e^y \quad \text{et} \quad u_2(x, y) = (2y - y^2 + 2)e^x \]
où $x, y \ge 0$.
\begin{itemize}
    \item Pour le joueur 1, maximiser $u_1$ par rapport à $x$ (en traitant $e^y$ comme constant positif) revient à maximiser $-x^2 + x + 1$.
    \[ \frac{\partial u_1}{\partial x} = (-2x + 1)e^y = 0 \implies x = 1/2 \]
    \item Pour le joueur 2, maximiser $u_2$ par rapport à $y$ revient à maximiser $2y - y^2 + 2$.
    \[ \frac{\partial u_2}{\partial y} = (2 - 2y)e^x = 0 \implies y = 1 \]
\end{itemize}
L'unique équilibre en stratégies dominantes est $(1/2, 1)$.
\end{exemple}

\begin{exemple}[Illustration du processus IESDS]
Considérons le jeu et la matrice suivants :
\[
\begin{pmatrix} & A & B & C \\ a & (1,-1) & (0,1) & (-1,0) \\ b & (2,3) & (2,2) & (0,1) \\ c & (1,4) & (0,3) & (1,2) \end{pmatrix}
\]
\textbf{Résolution par étapes :}
\begin{enumerate}
    \item \textbf{Étape 1 :} Pour le Joueur 1, l'action $a$ est strictement dominée par $b$ (car $1<2, 0<2, -1<0$). On élimine la ligne $a$.
    \item \textbf{Étape 2 :} Dans le jeu réduit (lignes $b, c$), le Joueur 2 compare ses colonnes.
    \begin{itemize}
        \item $B$ vs $A$ : $(2 < 3)$ si $b$, $(3 < 4)$ si $c$. Donc $B$ est dominée par $A$.
        \item $C$ vs $A$ : $(1 < 3)$ si $b$, $(2 < 4)$ si $c$. Donc $C$ est dominée par $A$.
    \end{itemize}
    On élimine les colonnes $B$ et $C$.
    \item \textbf{Étape 3 :} Dans le jeu réduit (colonne $A$ uniquement), le Joueur 1 choisit entre $b$ (gain 2) et $c$ (gain 1). L'action $c$ est dominée. On élimine $c$.
    \item \textbf{Solution :} Le seul profil survivant est $(b, A)$ avec gains $(2, 3)$.
\end{enumerate}
\end{exemple}

\section{Relation avec la Meilleure Réponse}

\begin{proposition}
Une action strictement dominée n'est \textbf{jamais} une meilleure réponse à une quelconque croyance sur les stratégies des adversaires.
\end{proposition}
\begin{remarque}
Lien fondamental : Cette proposition lie l'hypothèse comportementale de rationalité (ne pas jouer d'action dominée) au concept mathématique de meilleure réponse (point fixe, Nash).
\end{remarque}
\begin{exercice}[Analyse des jeux classiques]
\begin{itemize}
    \item \textbf{Dilemme du prisonnier :} Démontrez que "Confesser" est une stratégie dominante pour chaque joueur.
    \item \textbf{Jeu de coordination et Poule mouillée :} Montrez qu'aucun joueur ne possède de stratégie dominante (la meilleure action dépend de l'autre).
\end{itemize}
\end{exercice}

\begin{exemple}[Exam 2023 : IESDA Matrice]
\textbf{Question :}
Soit un jeu à deux joueurs défini par la matrice de gains suivante (le premier chiffre représente le gain du Joueur 1, le second celui du Joueur 2) :

\begin{center}
\begin{tabular}{c|cccc}
P1 \textbackslash P2 & W & X & Y & Z \\ \hline
A & (1, 0) & (1, 0) & (0, 0) & (-1, 0) \\
B & (4, 0) & (4, 0) & (1, 2) & (0, 3) \\
C & (2, 0) & (2, 0) & (0, 2) & (0, 3) \\
D & (3, 0) & (3, 0) & (3, 2) & (4, 4)
\end{tabular}
\end{center}

Déterminez l'équilibre de Nash en utilisant l'élimination itérée des stratégies strictement dominées.

\textbf{Correction :}
\begin{enumerate}
    \item \textbf{Étape 1 :} Pour le Joueur 1, l'action $A$ est strictement dominée par $B$ (car $4 > 1$, $4 > 1$, $1 > 0$, $0 > -1$). On élimine la ligne $A$.
    \item \textbf{Étape 2 :} Pour le Joueur 2, dans le jeu restant, on élimine l'action $X$ (dominée par d'autres colonnes).
    \item \textbf{Étape 3 :} Pour le Joueur 1, on élimine l'action $C$.
    \item \textbf{Étape 4 :} Pour le Joueur 2, on élimine l'action $W$.
    \item \textbf{Étape 5 :} Pour le Joueur 1, on élimine l'action $B$.
\end{enumerate}
\textbf{Résultat :} Le profil $(D, Z)$ est l'équilibre unique obtenu par élimination. On vérifie facilement que $(D, Z)$ est un équilibre de Nash.

\end{exemple}

\begin{exemple}[Preuve Théorique]
Soit $G$ un jeu où $u_i$ sont continues et $X_i$ compacts non vides.
\textbf{Question :} Prouver que l'ensemble des équilibres survivant à l'IESDA est compact et non vide.

\paragraph{Correction :}
Soit $X^0 = \prod X_i$ compact non vide.
À chaque étape $k$, on retire les stratégies strictement dominées. Si $u_i$ est continue, l'ensemble des stratégies dominées est un ouvert (ou une union d'ouverts), donc l'ensemble restant $X^k$ est fermé.
Comme $X^k \subset X^{k-1} \subset \dots \subset X^0$, $X^k$ est une suite décroissante de compacts.
Si l'on suppose que l'élimination est bien définie (jamais vide à chaque étape finie grâce au théorème du maximum), alors l'intersection infinie $X^\infty = \bigcap_{k \ge 0} X^k$ est une intersection décroissante de compacts non vides dans un espace séparé.
D'après la propriété de Borel-Lebesgue (ou théorème de Cantor), cette intersection est \textbf{non vide et compacte}.
\end{exemple}

\begin{exemple}[Vérification Classique]
Existe-t-il des stratégies dominantes dans :
\begin{itemize}
    \item Le jeu de la Coordination ? \textbf{Non.} (Dépend de l'autre).
    \item Le Chicken Game (Poule mouillée) ? \textbf{Non.} (Si l'autre fonce, je dévie ; s'il dévie, je fonce).
    \item Le Dilemme du Prisonnier ? \textbf{Oui.} "Dénoncer" est une stratégie strictement dominante pour les deux joueurs, quelle que soit l'action de l'autre.
\end{itemize}
\end{exemple}
